As quantum computing technologies continue to advance and quantum communication links become increasingly relevant for interconnecting remote quantum systems, the reliable distribution of quantum states over optical fibers is becoming a key technological requirement.

Polarization-entangled photon pairs are well suited to fiber links because polarization can be manipulated and measured with standard optical components. Real fibers, however, introduce coherent transformations and decoherence mechanisms that modify the observed correlations and may degrade the distributed entanglement. This thesis develops a progressive framework connecting these physical effects to numerical predictions and experimentally accessible quantities.

The modeling starts from an experimentally motivated description of the residual polarization evolution remaining after partial compensation. A fixed relative phase is first introduced to distinguish coherent changes in measured correlations from actual entanglement degradation. The model is then extended to unresolved phase fluctuations, summarized by the coherence parameter \( \mu=\langle e^{i\theta}\rangle \), whose magnitude determines the residual concurrence within this model through \( C=|\mu| \). Independent depolarizing effects in the two fiber arms are subsequently introduced through the correlation-survival factor \( \eta=(1-p_A)(1-p_B) \) and combined with phase fluctuations in a common density-operator description. Finally, the effective models are connected to polarization-mode dispersion, where finite-bandwidth photons undergo frequency-dependent birefringent transformations and polarization mixedness emerges after tracing over unresolved spectral degrees of freedom.

The framework is implemented in a standalone modular MATLAB library reproducing this modeling hierarchy numerically. The simulator provides a common pipeline for state preparation, channel composition, statistical averaging, polarization-correlation analysis, and evaluation of purity, fidelity, concurrence, and Bell-inequality parameters. Analytical predictions are compared with numerical results to validate the implementation and identify the different signatures of coherent evolution, statistical dephasing, depolarization, and PMD.

The theoretical and computational work is complemented by the development of a two-arm laboratory platform for polarization-entanglement distribution based on single-photon detection and coincidence measurements together with controllable polarization transformations. The setup provides the experimental counterpart of the developed framework and a basis for a future entanglement-based BBM92 quantum-key-distribution implementation. Overall, the thesis establishes a continuous path from fiber-channel modeling, through numerical simulation, to the experimental characterization of entanglement distribution in optical fibers.